% This LaTeX file was created by Metamath on 6-Mar-2024 11:48 AM.

\dfc
%\begin{restatable}{metadefinition}{dfc}~\otherTitle{df-c}~ Define the set of complex numbers.  \begin{align}
%\vdash \mathbb{C} = ( \mathcal{R} \times \mathcal{R} )\label{eq:df-c}\tag{df-c}
%\end{align}
%\end{restatable}

\dfr
%\begin{restatable}{metadefinition}{dfr}~\otherTitle{df-r}~ Define the set of real numbers. 
%\begin{align}
%\vdash \mathbb{R} = ( \mathcal{R} \times \{ 0_\mathcal{R} \}
%    )\label{eq:df-r}\tag{df-r}
%\end{align}
%\end{restatable}


\begin{restatable}{axiom}{axresscn}~\otherTitle{ax-resscn}~ The real numbers are a subset of the
complex numbers. 
\begin{align}
\vdash \mathbb{R} \subseteq \mathbb{C}\label{eq:ax-resscn}\tag{ax-resscn}
\end{align}
\end{restatable}

\begin{restatable}{lemma}{0cn}~\otherTitle{0cn}~ Zero is a complex number. 
\begin{align}
\vdash 0 \in \mathbb{C}\label{eq:0cn}\tag{0cn}
\end{align}
\end{restatable}

\begin{restatable}{lemma}{addcan}~\otherTitle{addcan}~ Cancellation law for addition.  T\ ~
\begin{align}
\vdash ( ( \mc{A} \in \mathbb{C} \wedge \mc{B} \in \mathbb{C} \wedge \mc{C} \in
    \mathbb{C} ) \rightarrow ( ( \mc{A} + \mc{B} ) = ( \mc{A} + \mc{C} )
    \leftrightarrow \mc{B} = \mc{C} ) )\label{eq:addcan}\tag{addcan}
\end{align}
\end{restatable}

\begin{restatable}{lemma}{add12}~\otherTitle{add12}~ Commutative/associative law that swaps the
first two terms in a triple
     sum.  ~
\begin{align}
\vdash ( ( \mc{A} \in \mathbb{C} \wedge \mc{B} \in \mathbb{C} \wedge \mc{C} \in
    \mathbb{C} ) \rightarrow ( \mc{A} + ( \mc{B} + \mc{C} ) ) = ( \mc{B} + (
    \mc{A} + \mc{C} ) ) )\label{eq:add12}\tag{add12}
\end{align}
\end{restatable}

\begin{restatable}{lemma}{addcl}~\otherTitle{addcl}~ Addition of Complex Numbers is Complex.~
\begin{align}
\vdash ( ( \mc{A} \in \mathbb{C} \wedge \mc{B} \in \mathbb{C} ) \rightarrow (
    \mc{A} + \mc{B} ) \in \mathbb{C} )\label{eq:addcl}\tag{addcl}
\end{align}
\end{restatable}

\begin{restatable}{lemma}{readdcl}~\otherTitle{readdcl}~ Addition of Real Numbers is Real,     ~
\begin{align}
\vdash ( ( \mc{A} \in \mathbb{R} \wedge \mc{B} \in \mathbb{R} ) \rightarrow (
    \mc{A} + \mc{B} ) \in \mathbb{R} )\label{eq:readdcl}\tag{readdcl}
\end{align}
\end{restatable}

\begin{restatable}{lemma}{mulcl}~\otherTitle{mulcl}~ Multiplication of Complex Numbers is Complex.     ~
\begin{align}
\vdash ( ( \mc{A} \in \mathbb{C} \wedge \mc{B} \in \mathbb{C} ) \rightarrow (
    \mc{A} \times \mc{B} ) \in \mathbb{C} )\label{eq:mulcl}\tag{mulcl}
\end{align}
\end{restatable}

\begin{restatable}{lemma}{remulcl}~\otherTitle{remulcl}~ Multitplication  of Real Numbers is Real,    ~
\begin{align}
\vdash ( ( \mc{A} \in \mathbb{R} \wedge \mc{B} \in \mathbb{R} ) \rightarrow (
    \mc{A} \times \mc{B} ) \in \mathbb{R} )\label{eq:remulcl}\tag{remulcl}
\end{align}
\end{restatable}

\begin{restatable}{lemma}{mulcom}~\otherTitle{mulcom}~ Multiplication Commutes.     ~
\begin{align}
\vdash ( ( \mc{A} \in \mathbb{C} \wedge \mc{B} \in \mathbb{C} ) \rightarrow (
    \mc{A} \times \mc{B} ) = ( \mc{B} \times \mc{A} )
    )\label{eq:mulcom}\tag{mulcom}
\end{align}
\end{restatable}

\begin{restatable}{lemma}{mulass}~\otherTitle{mulass}~ 
Multiplication is Associative.
\begin{align}
\vdash ( ( \mc{A} \in \mathbb{C} \wedge \mc{B} \in \mathbb{C} \wedge \mc{C} \in
    \mathbb{C} ) \rightarrow ( ( \mc{A} \times \mc{B} ) \times \mc{C} ) = (
    \mc{A} \times ( \mc{B} \times \mc{C} ) ) )\label{eq:mulass}\tag{mulass}
\end{align}
\end{restatable}

\begin{restatable}{lemma}{recn}~\otherTitle{recn}~ A real number is a complex number. 
~
\begin{align}
\vdash ( \mc{A} \in \mathbb{R} \rightarrow \mc{A} \in \mathbb{C}
    )\label{eq:recn}\tag{recn}
\end{align}
\end{restatable}

\begin{restatable}{lemma}{0re}~\otherTitle{0re}~ The number 0 is real.   ~
\begin{align}
\vdash 0 \in \mathbb{R}\label{eq:0re}\tag{0re}
\end{align}
\end{restatable}

\begin{restatable}{lemma}{addcli}~\otherTitle{addcli}~ Closure law for addition.  ~
\begin{align}
\vdash \mc{A} \in \mathbb{C}\quad\&\quad \vdash \mc{B} \in
    \mathbb{C}\quad\Rightarrow\quad \vdash ( \mc{A} + \mc{B} ) \in
    \mathbb{C}\label{eq:addcli}\tag{addcli}
\end{align}
\end{restatable}

\begin{restatable}{lemma}{mulcli}~\otherTitle{mulcli}~ Closure law for multiplication. 
~
\begin{align}
\vdash \mc{A} \in \mathbb{C}\quad\&\quad \vdash \mc{B} \in
    \mathbb{C}\quad\Rightarrow\quad \vdash ( \mc{A} \times \mc{B} ) \in
    \mathbb{C}\label{eq:mulcli}\tag{mulcli}
\end{align}
\end{restatable}

\begin{restatable}{lemma}{mulcomi}~\otherTitle{mulcomi}~ Commutative law for multiplication. 
~
\begin{align}
\vdash \mc{A} \in \mathbb{C}\quad\&\quad \vdash \mc{B} \in
    \mathbb{C}\quad\Rightarrow\quad \vdash ( \mc{A} \times \mc{B} ) = ( \mc{B}
    \times \mc{A} )\label{eq:mulcomi}\tag{mulcomi}
\end{align}
\end{restatable}

\begin{restatable}{lemma}{addassi}~\otherTitle{addassi}~ Associative law for addition. 
\begin{align}
\vdash \mc{A} \in \mathbb{C}\quad\&\quad \vdash \mc{B} \in
    \mathbb{C}\quad\&\quad \vdash \mc{C} \in \mathbb{C}\quad\Rightarrow\quad
    \vdash ( ( \mc{A} + \mc{B} ) + \mc{C} ) = ( \mc{A} + ( \mc{B} + \mc{C} )
    )\label{eq:addassi}\tag{addassi}
\end{align}
\end{restatable}

\begin{restatable}{lemma}{adddii}~\otherTitle{adddii}~ Distributive law (left-distributivity). 
~
\begin{align}
\vdash \mc{A} \in \mathbb{C}\quad\&\quad \vdash \mc{B} \in
    \mathbb{C}\quad\&\quad \vdash \mc{C} \in \mathbb{C}\quad\Rightarrow\quad
    \vdash ( \mc{A} \times ( \mc{B} + \mc{C} ) ) = ( ( \mc{A} \times \mc{B} ) + (
    \mc{A} \times \mc{C} ) )\label{eq:adddii}\tag{adddii}
\end{align}
\end{restatable}

\begin{restatable}{lemma}{recni}~\otherTitle{recni}~ A real number is a complex number. 
~
\begin{align}
\vdash \mc{A} \in \mathbb{R}\quad\Rightarrow\quad \vdash \mc{A} \in
    \mathbb{C}\label{eq:recni}\tag{recni}
\end{align}
\end{restatable}

\begin{restatable}{lemma}{remulcli}~\otherTitle{remulcli}~ Closure law for multiplication of reals. 
~
\begin{align}
\vdash \mc{A} \in \mathbb{R}\quad\&\quad \vdash \mc{B} \in
    \mathbb{R}\quad\Rightarrow\quad \vdash ( \mc{A} \times \mc{B} ) \in
    \mathbb{R}\label{eq:remulcli}\tag{remulcli}
\end{align}
\end{restatable}

\begin{restatable}{metadefinition}{dfltxr}~\otherTitle{df-ltxr}~ Define 'less than' on the set of
extended reals.        ~
\begin{align}
\vdash < = ( \{ \langle \msc{x} , \msc{y} \rangle~\vert~( \msc{x} \in \mathbb{R}
    \wedge \msc{y} \in \mathbb{R} \wedge \msc{x} <_\mathbb{R} \msc{y} ) \} \cup
    ( ( ( \mathbb{R} \cup \{ -\infty \} ) \times \{ +\infty \} ) \cup ( \{
    -\infty \} \times \mathbb{R} ) ) )\label{eq:df-ltxr}\tag{df-ltxr}
\end{align}
\end{restatable}

\begin{restatable}{lemma}{ressxr}~\otherTitle{ressxr}~ The standard reals are a subset of the
extended reals.  ~
\begin{align}
\vdash \mathbb{R} \subseteq \mathbb{R}^{\ast}\label{eq:ressxr}\tag{ressxr}
\end{align}
\end{restatable}

\begin{restatable}{lemma}{lttr}~\otherTitle{lttr}~ Less than is transitive.
\begin{align}
\vdash ( ( \mc{A} \in \mathbb{R} \wedge \mc{B} \in \mathbb{R} \wedge \mc{C} \in
    \mathbb{R} ) \rightarrow ( ( \mc{A} < \mc{B} \wedge \mc{B} < \mc{C} )
    \rightarrow \mc{A} < \mc{C} ) )\label{eq:lttr}\tag{lttr}
\end{align}
\end{restatable}

\begin{restatable}{lemma}{leloe}~\otherTitle{leloe}~ 'Less than or equal to' expressed in terms of
'less than' or 'equals'.
     ~
\begin{align}
\vdash ( ( \mc{A} \in \mathbb{R} \wedge \mc{B} \in \mathbb{R} ) \rightarrow (
    \mc{A} \le \mc{B} \leftrightarrow ( \mc{A} < \mc{B} \vee \mc{A} = \mc{B} )
    ) )\label{eq:leloe}\tag{leloe}
\end{align}
\end{restatable}

\begin{restatable}{lemma}{ltle}~\otherTitle{ltle}~ 'Less than' implies 'less than or equal to'. 
~
\begin{align}
\vdash ( ( \mc{A} \in \mathbb{R} \wedge \mc{B} \in \mathbb{R} ) \rightarrow (
    \mc{A} < \mc{B} \rightarrow \mc{A} \le \mc{B} ) )\label{eq:ltle}\tag{ltle}
\end{align}
\end{restatable}

\begin{restatable}{lemma}{ltnr}~\otherTitle{ltnr}~ 'Less than' is irreflexive.  ~
\begin{align}
\vdash ( \mc{A} \in \mathbb{R} \rightarrow \lnot \mc{A} < \mc{A}
    )\label{eq:ltnr}\tag{ltnr}
\end{align}
\end{restatable}

\begin{restatable}{lemma}{leid}~\otherTitle{leid}~ 'Less than or equal to' is reflexive. 
~
\begin{align}
\vdash ( \mc{A} \in \mathbb{R} \rightarrow \mc{A} \le \mc{A}
    )\label{eq:leid}\tag{leid}
\end{align}
\end{restatable}

\begin{restatable}{lemma}{lenlti}~\otherTitle{lenlti}~ 'Less than or equal to' in terms of 'less
than'.  ~
\begin{align}
\vdash \mc{A} \in \mathbb{R}\quad\&\quad \vdash \mc{B} \in
    \mathbb{R}\quad\Rightarrow\quad \vdash ( \mc{A} \le \mc{B} \leftrightarrow
    \lnot \mc{B} < \mc{A} )\label{eq:lenlti}\tag{lenlti}
\end{align}
\end{restatable}

\begin{restatable}{lemma}{ltnei}~\otherTitle{ltnei}~ 'Less than' implies not equal. 
\begin{align}
\vdash \mc{A} \in \mathbb{R}\quad\&\quad \vdash \mc{B} \in
    \mathbb{R}\quad\Rightarrow\quad \vdash ( \mc{A} < \mc{B} \rightarrow \mc{B}
    \ne \mc{A} )\label{eq:ltnei}\tag{ltnei}
\end{align}
\end{restatable}

\begin{restatable}{lemma}{letri}~\otherTitle{letri}~ 'Less than or equal to' is transitive. 
~
\begin{align}
\vdash \mc{A} \in \mathbb{R}\quad\&\quad \vdash \mc{B} \in
    \mathbb{R}\quad\&\quad \vdash \mc{C} \in \mathbb{R}\quad\Rightarrow\quad
    \vdash ( ( \mc{A} \le \mc{B} \wedge \mc{B} \le \mc{C} ) \rightarrow \mc{A}
    \le \mc{C} )\label{eq:letri}\tag{letri}
\end{align}
\end{restatable}

\begin{restatable}{lemma}{addid1}~\otherTitle{addid1}~ $0$ is an additive identity.  
       ~
\begin{align}
\vdash ( \mc{A} \in \mathbb{C} \rightarrow ( \mc{A} + 0 ) = \mc{A}
    )\label{eq:addid1}\tag{addid1}
\end{align}
\end{restatable}

\begin{restatable}{lemma}{mul12i}~\otherTitle{mul12i}~ Commutative/associative law that swaps the
first two factors in a triple
       product.  ~  ~
\begin{align}
\vdash \mc{A} \in \mathbb{C}\quad\&\quad \vdash \mc{B} \in
    \mathbb{C}\quad\&\quad \vdash \mc{C} \in \mathbb{C}\quad\Rightarrow\quad
    \vdash ( \mc{A} \times ( \mc{B} \times \mc{C} ) ) = ( \mc{B} \times ( \mc{A}
    \times \mc{C} ) )\label{eq:mul12i}\tag{mul12i}
\end{align}
\end{restatable}

\begin{restatable}{metadefinition}{dfneg}~\otherTitle{df-neg}~ Define the negative of a number (unary
minus).  
\begin{align}
\vdash \textrm{-} \mc{A} = ( 0 \minus \mc{A} )\label{eq:df-neg}\tag{df-neg}
\end{align}
\end{restatable}

\begin{restatable}{lemma}{negeqi}~\otherTitle{negeqi}~ Equality inference for negatives. 
~
\begin{align}
\vdash \mc{A} = \mc{B}\quad\Rightarrow\quad \vdash \textrm{-} \mc{A} =
    \textrm{-} \mc{B}\label{eq:negeqi}\tag{negeqi}
\end{align}
\end{restatable}

\begin{restatable}{lemma}{negcl}~\otherTitle{negcl}~ Closure law for negative.  ~
\begin{align}
\vdash ( \mc{A} \in \mathbb{C} \rightarrow \textrm{-} \mc{A} \in \mathbb{C}
    )\label{eq:negcl}\tag{negcl}
\end{align}
\end{restatable}

\begin{restatable}{lemma}{subcl}~\otherTitle{subcl}~ Closure law for subtraction.  ~
       ~
\begin{align}
\vdash ( ( \mc{A} \in \mathbb{C} \wedge \mc{B} \in \mathbb{C} ) \rightarrow (
    \mc{A} \minus \mc{B} ) \in \mathbb{C} )\label{eq:subcl}\tag{subcl}
\end{align}
\end{restatable}

\begin{restatable}{lemma}{mul01}~\otherTitle{mul01}~ Multiplication by $0$.    ~
\begin{align}
\vdash ( \mc{A} \in \mathbb{C} \rightarrow ( \mc{A} \times 0 ) = 0
    )\label{eq:mul01}\tag{mul01}
\end{align}
\end{restatable}

\begin{restatable}{lemma}{pncan}~\otherTitle{pncan}~ Cancellation law for subtraction. 
~
     ~
\begin{align}
\vdash ( ( \mc{A} \in \mathbb{C} \wedge \mc{B} \in \mathbb{C} ) \rightarrow ( (
    \mc{A} + \mc{B} ) \minus \mc{B} ) = \mc{A} )\label{eq:pncan}\tag{pncan}
\end{align}
\end{restatable}

\begin{restatable}{lemma}{pncan2}~\otherTitle{pncan2}~ Cancellation law for subtraction. 
~
\begin{align}
\vdash ( ( \mc{A} \in \mathbb{C} \wedge \mc{B} \in \mathbb{C} ) \rightarrow ( (
    \mc{A} + \mc{B} ) \minus \mc{A} ) = \mc{B} )\label{eq:pncan2}\tag{pncan2}
\end{align}
\end{restatable}

\begin{restatable}{lemma}{npcan}~\otherTitle{npcan}~ Cancellation law for subtraction. 
~
     ~
\begin{align}
\vdash ( ( \mc{A} \in \mathbb{C} \wedge \mc{B} \in \mathbb{C} ) \rightarrow ( (
    \mc{A} \minus \mc{B} ) + \mc{B} ) = \mc{A} )\label{eq:npcan}\tag{npcan}
\end{align}
\end{restatable}

\begin{restatable}{lemma}{subid}~\otherTitle{subid}~ Subtraction of a number from itself. 
~
     ~
\begin{align}
\vdash ( \mc{A} \in \mathbb{C} \rightarrow ( \mc{A} \minus \mc{A} ) = 0
    )\label{eq:subid}\tag{subid}
\end{align}
\end{restatable}

\begin{restatable}{lemma}{subid1}~\otherTitle{subid1}~ Identity law for subtraction.  
\begin{align}
\vdash ( \mc{A} \in \mathbb{C} \rightarrow ( \mc{A} \minus 0 ) = \mc{A}
    )\label{eq:subid1}\tag{subid1}
\end{align}
\end{restatable}

\begin{restatable}{lemma}{subcan2}~\otherTitle{subcan2}~ Cancellation law for subtraction. 
~
\begin{align}
\vdash ( ( \mc{A} \in \mathbb{C} \wedge \mc{B} \in \mathbb{C} \wedge \mc{C} \in
    \mathbb{C} ) \rightarrow ( ( \mc{A} \minus \mc{C} ) = ( \mc{B} \minus
    \mc{C} ) \leftrightarrow \mc{A} = \mc{B} ) )\label{eq:subcan2}\tag{subcan2}
\end{align}
\end{restatable}

\begin{restatable}{lemma}{subsub3}~\otherTitle{subsub3}~ Law for double subtraction.  
\begin{align}
\vdash ( ( \mc{A} \in \mathbb{C} \wedge \mc{B} \in \mathbb{C} \wedge \mc{C} \in
    \mathbb{C} ) \rightarrow ( \mc{A} \minus ( \mc{B} \minus \mc{C} ) ) = ( (
    \mc{A} + \mc{C} ) \minus \mc{B} ) )\label{eq:subsub3}\tag{subsub3}
\end{align}
\end{restatable}

\begin{restatable}{lemma}{subsub4}~\otherTitle{subsub4}~ Law for double subtraction.  
\begin{align}
\vdash ( ( \mc{A} \in \mathbb{C} \wedge \mc{B} \in \mathbb{C} \wedge \mc{C} \in
    \mathbb{C} ) \rightarrow ( ( \mc{A} \minus \mc{B} ) \minus \mc{C} ) = (
    \mc{A} \minus ( \mc{B} + \mc{C} ) ) )\label{eq:subsub4}\tag{subsub4}
\end{align}
\end{restatable}

\begin{restatable}{lemma}{addsub4}~\otherTitle{addsub4}~ Rearrangement of 4 terms in a mixed
addition and subtraction.
     ~
\begin{align}
\vdash ( ( ( \mc{A} \in \mathbb{C} \wedge \mc{B} \in \mathbb{C} ) \wedge (
    \mc{C} \in \mathbb{C} \wedge \mc{D} \in \mathbb{C} ) ) \rightarrow ( (
    \mc{A} + \mc{B} ) \minus ( \mc{C} + \mc{D} ) ) = ( ( \mc{A} \minus \mc{C} )
    + ( \mc{B} \minus \mc{D} ) ) )\label{eq:addsub4}\tag{addsub4}
\end{align}
\end{restatable}

\begin{restatable}{lemma}{subadd4}~\otherTitle{subadd4}~ Rearrangement of 4 terms in a mixed
addition and subtraction.
     ~
\begin{align}
\vdash ( ( ( \mc{A} \in \mathbb{C} \wedge \mc{B} \in \mathbb{C} ) \wedge (
    \mc{C} \in \mathbb{C} \wedge \mc{D} \in \mathbb{C} ) ) \rightarrow ( (
    \mc{A} \minus \mc{B} ) \minus ( \mc{C} \minus \mc{D} ) ) = ( ( \mc{A} +
    \mc{D} ) \minus ( \mc{B} + \mc{C} ) ) )\label{eq:subadd4}\tag{subadd4}
\end{align}
\end{restatable}

\begin{restatable}{lemma}{negid}~\otherTitle{negid}~ Addition of a number and its negative. 
~
\begin{align}
\vdash ( \mc{A} \in \mathbb{C} \rightarrow ( \mc{A} + \textrm{-} \mc{A} ) = 0
    )\label{eq:negid}\tag{negid}
\end{align}
\end{restatable}

\begin{restatable}{lemma}{negsub}~\otherTitle{negsub}~ Relationship between subtraction and
negative.    ~
\begin{align}
\vdash ( ( \mc{A} \in \mathbb{C} \wedge \mc{B} \in \mathbb{C} ) \rightarrow (
    \mc{A} + \textrm{-} \mc{B} ) = ( \mc{A} \minus \mc{B} )
    )\label{eq:negsub}\tag{negsub}
\end{align}
\end{restatable}

\begin{restatable}{lemma}{negdi}~\otherTitle{negdi}~ Distribution of negative over addition. 
~
     ~
\begin{align}
\vdash ( ( \mc{A} \in \mathbb{C} \wedge \mc{B} \in \mathbb{C} ) \rightarrow
    \textrm{-} ( \mc{A} + \mc{B} ) = ( \textrm{-} \mc{A} + \textrm{-} \mc{B} )
    )\label{eq:negdi}\tag{negdi}
\end{align}
\end{restatable}

\begin{restatable}{lemma}{negsubdi2}~\otherTitle{negsubdi2}~ Distribution of negative over
subtraction.  ~
\begin{align}
\vdash ( ( \mc{A} \in \mathbb{C} \wedge \mc{B} \in \mathbb{C} ) \rightarrow
    \textrm{-} ( \mc{A} \minus \mc{B} ) = ( \mc{B} \minus \mc{A} )
    )\label{eq:negsubdi2}\tag{negsubdi2}
\end{align}
\end{restatable}

\begin{restatable}{lemma}{negcli}~\otherTitle{negcli}~ Closure law for negative.  ~
\begin{align}
\vdash \mc{A} \in \mathbb{C}\quad\Rightarrow\quad \vdash \textrm{-} \mc{A} \in
    \mathbb{C}\label{eq:negcli}\tag{negcli}
\end{align}
\end{restatable}

\begin{restatable}{lemma}{pncan3i}~\otherTitle{pncan3i}~ Subtraction and addition of equals. 
~
\begin{align}
\vdash \mc{A} \in \mathbb{C}\quad\&\quad \vdash \mc{B} \in
    \mathbb{C}\quad\Rightarrow\quad \vdash ( \mc{A} + ( \mc{B} \minus \mc{A} )
    ) = \mc{B}\label{eq:pncan3i}\tag{pncan3i}
\end{align}
\end{restatable}

\begin{restatable}{lemma}{negcld}~\otherTitle{negcld}~ Closure law for negative.  ~
\begin{align}
\vdash ( \mw{\varphi} \rightarrow \mc{A} \in \mathbb{C} )\quad\Rightarrow\quad
    \vdash ( \mw{\varphi} \rightarrow \textrm{-} \mc{A} \in \mathbb{C}
    )\label{eq:negcld}\tag{negcld}
\end{align}
\end{restatable}

\begin{restatable}{lemma}{nncand}~\otherTitle{nncand}~ Cancellation law for subtraction. 
~
\begin{align}
\vdash ( \mw{\varphi} \rightarrow \mc{A} \in \mathbb{C} )\quad\&\quad \vdash (
    \mw{\varphi} \rightarrow \mc{B} \in \mathbb{C} )\quad\Rightarrow\quad
    \vdash ( \mw{\varphi} \rightarrow ( \mc{A} \minus ( \mc{A} \minus \mc{B} )
    ) = \mc{B} )\label{eq:nncand}\tag{nncand}
\end{align}
\end{restatable}

\begin{restatable}{lemma}{subnegd}~\otherTitle{subnegd}~ Relationship between subtraction and
negative.  ~
\begin{align}
\vdash ( \mw{\varphi} \rightarrow \mc{A} \in \mathbb{C} )\quad\&\quad \vdash (
    \mw{\varphi} \rightarrow \mc{B} \in \mathbb{C} )\quad\Rightarrow\quad
    \vdash ( \mw{\varphi} \rightarrow ( \mc{A} \minus \textrm{-} \mc{B} ) = (
    \mc{A} + \mc{B} ) )\label{eq:subnegd}\tag{subnegd}
\end{align}
\end{restatable}

\begin{restatable}{lemma}{mulcan2}~\otherTitle{mulcan2}~ Cancellation law for multiplication. 
~
     ~
\begin{align}
\vdash ( ( \mc{A} \in \mathbb{C} \wedge \mc{B} \in \mathbb{C} \wedge ( \mc{C}
    \in \mathbb{C} \wedge \mc{C} \ne 0 ) ) \rightarrow ( ( \mc{A} \times \mc{C}
    ) = ( \mc{B} \times \mc{C} ) \leftrightarrow \mc{A} = \mc{B} )
    )\label{eq:mulcan2}\tag{mulcan2}
\end{align}
\end{restatable}

\begin{restatable}{lemma}{mulle0b}~\otherTitle{mulle0b}~ A condition for multiplication to be
nonpositive.  ~
\begin{align}
\vdash ( ( \mc{A} \in \mathbb{R} \wedge \mc{B} \in \mathbb{R} ) \rightarrow ( (
    \mc{A} \times \mc{B} ) \le 0 \leftrightarrow ( ( \mc{A} \le 0 \wedge 0 \le
    \mc{B} ) \vee ( 0 \le \mc{A} \wedge \mc{B} \le 0 ) ) )
    )\label{eq:mulle0b}\tag{mulle0b}
\end{align}
\end{restatable}

\begin{restatable}{lemma}{0z}~\otherTitle{0z}~ Zero is an integer.  ~
\begin{align}
\vdash 0 \in \mathbb{Z}\label{eq:0z}\tag{0z}
\end{align}
\end{restatable}

\begin{restatable}{lemma}{subsub4}~\otherTitle{subsub4}~ Law for double subtraction.  
\begin{align}
\vdash ( ( \mc{A} \in \mathbb{C} \wedge \mc{B} \in \mathbb{C} \wedge \mc{C} \in
    \mathbb{C} ) \rightarrow ( ( \mc{A} \minus \mc{B} ) \minus \mc{C} ) = (
    \mc{A} \minus ( \mc{B} + \mc{C} ) ) )\label{eq:subsub4}\tag{subsub4}
\end{align}
\end{restatable}

\begin{restatable}{lemma}{readdcli}~\otherTitle{readdcli}~ Closure law for addition of reals. 
~
\begin{align}
\vdash \mc{A} \in \mathbb{R}\quad\&\quad \vdash \mc{B} \in
    \mathbb{R}\quad\Rightarrow\quad \vdash ( \mc{A} + \mc{B} ) \in
    \mathbb{R}\label{eq:readdcli}\tag{readdcli}
\end{align}
\end{restatable}

\begin{restatable}{lemma}{letri3}~\otherTitle{letri3}~ Trichotomy law.  ~
\begin{align}
\vdash ( ( \mc{A} \in \mathbb{R} \wedge \mc{B} \in \mathbb{R} ) \rightarrow (
    \mc{A} = \mc{B} \leftrightarrow ( \mc{A} \le \mc{B} \wedge \mc{B} \le
    \mc{A} ) ) )\label{eq:letri3}\tag{letri3}
\end{align}
\end{restatable}

\begin{restatable}{lemma}{leadd2i}~\otherTitle{leadd2i}~ Addition to both sides of 'less than or
equal to'.  ~
\begin{align}
\vdash \mc{A} \in \mathbb{R}\quad\&\quad \vdash \mc{B} \in
    \mathbb{R}\quad\&\quad \vdash \mc{C} \in \mathbb{R}\quad\Rightarrow\quad
    \vdash ( \mc{A} \le \mc{B} \leftrightarrow ( \mc{C} + \mc{A} ) \le ( \mc{C}
    + \mc{B} ) )\label{eq:leadd2i}\tag{leadd2i}
\end{align}
\end{restatable}

\dfnn
%\begin{restatable}{metadefinition}{dfnn}~\otherTitle{df-nn}~ Define the set of positive integers.  
%\begin{align}
%\vdash \mathbb{N} = ( \,\mathrm{rec} ( ( \msc{x} \in \,\mathrm{V} \mapsto (
%    \msc{x} + 1 ) ) , 1 ) `` \,\mathrm{\omega} )\label{eq:df-nn}\tag{df-nn}
%\end{align}
%\end{restatable}

\dfno
%\begin{restatable}{metadefinition}{dfno}~\otherTitle{df-n0}~ Define the set of nonnegative integers. 
%~
%\begin{align}
%\vdash \mathbb{N}_0 = ( \mathbb{N} \cup \{ 0 \} )\label{eq:df-n0}\tag{df-n0}
%\end{align}
%\end{restatable}

\dfz
%\begin{restatable}{metadefinition}{dfz}~\otherTitle{df-z}~ Define the set of integers, which are the
%positive and negative integers
%     together with zero.  
%\begin{align}
%\vdash \mathbb{Z} = \{ \msc{n} \in \mathbb{R}~\vert~( \msc{n} = 0 \vee \msc{n} \in
%    \mathbb{N} \vee \textrm{-} \msc{n} \in \mathbb{N} )
%    \}\label{eq:df-z}\tag{df-z}
%\end{align}
%\end{restatable}

\dfq

\begin{restatable}{lemma}{zre}~\otherTitle{zre}~ An integer is a real.  ~
\begin{align}
\vdash ( N \in \mathbb{Z} \rightarrow N \in \mathbb{R} )\label{eq:zre}\tag{zre}
\end{align}
\end{restatable}

\begin{restatable}{lemma}{zcn}~\otherTitle{zcn}~ An integer is a complex number.  
\begin{align}
\vdash ( N \in \mathbb{Z} \rightarrow N \in \mathbb{C} )\label{eq:zcn}\tag{zcn}
\end{align}
\end{restatable}

\begin{restatable}{lemma}{zsscn}~\otherTitle{zsscn}~ The integers are a subset of the complex
numbers.  ~
\begin{align}
\vdash \mathbb{Z} \subseteq \mathbb{C}\label{eq:zsscn}\tag{zsscn}
\end{align}
\end{restatable}

\begin{restatable}{lemma}{zaddcl}~\otherTitle{zaddcl}~ Closure of addition of integers. 
\begin{align}
\vdash ( ( M \in \mathbb{Z} \wedge N \in \mathbb{Z} ) \rightarrow ( M + N ) \in
    \mathbb{Z} )\label{eq:zaddcl}\tag{zaddcl}
\end{align}
\end{restatable}


\begin{restatable}{lemma}{elioo1}~\otherTitle{elioo1}~ Membership in an open interval of extended
reals.  ~  ~
\begin{align}
\vdash ( ( \mc{A} \in \mathbb{R}^{\ast} \wedge \mc{B} \in \mathbb{R}^{\ast} )
    \rightarrow ( \mc{C} \in ( \mc{A} (,) \mc{B} ) \leftrightarrow ( \mc{C} \in
    \mathbb{R}^{\ast} \wedge \mc{A} < \mc{C} \wedge \mc{C} < \mc{B} ) )
    )\label{eq:elioo1}\tag{elioo1}
\end{align}
\end{restatable}

\begin{restatable}{lemma}{elioc1}~\otherTitle{elioc1}~ Membership in an open-below, closed-above
interval of extended reals.
       ~  ~
\begin{align}
\vdash ( ( \mc{A} \in \mathbb{R}^{\ast} \wedge \mc{B} \in \mathbb{R}^{\ast} )
    \rightarrow ( \mc{C} \in ( \mc{A} (,] \mc{B} ) \leftrightarrow ( \mc{C} \in
    \mathbb{R}^{\ast} \wedge \mc{A} < \mc{C} \wedge \mc{C} \le \mc{B} ) )
    )\label{eq:elioc1}\tag{elioc1}
\end{align}
\end{restatable}

\begin{restatable}{lemma}{elico1}~\otherTitle{elico1}~ Membership in a closed-below, open-above
interval of extended reals.
       ~  ~
\begin{align}
\vdash ( ( \mc{A} \in \mathbb{R}^{\ast} \wedge \mc{B} \in \mathbb{R}^{\ast} )
    \rightarrow ( \mc{C} \in ( \mc{A} [,) \mc{B} ) \leftrightarrow ( \mc{C} \in
    \mathbb{R}^{\ast} \wedge \mc{A} \le \mc{C} \wedge \mc{C} < \mc{B} ) )
    )\label{eq:elico1}\tag{elico1}
\end{align}
\end{restatable}

\begin{restatable}{lemma}{elicc1}~\otherTitle{elicc1}~ Membership in a closed interval of extended
reals.  ~  ~
\begin{align}
\vdash ( ( \mc{A} \in \mathbb{R}^{\ast} \wedge \mc{B} \in \mathbb{R}^{\ast} )
    \rightarrow ( \mc{C} \in ( \mc{A} [,] \mc{B} ) \leftrightarrow ( \mc{C} \in
    \mathbb{R}^{\ast} \wedge \mc{A} \le \mc{C} \wedge \mc{C} \le \mc{B} ) )
    )\label{eq:elicc1}\tag{elicc1}
\end{align}
\end{restatable}

